\section{Verificación del reloj de Lamport}


\subsection{Método}
El script \texttt{analysis/verificar\_lamport.py} lee los CSV de eventos reales
de \texttt{data/} y analiza cada corrida por separado. Cada envío
(\texttt{send}) se empareja con su recepción (\texttt{receive}) cuando
coinciden el identificador de corrida (\texttt{run\_id}), el identificador de
mensaje (\texttt{message\_id}) y el tipo (\texttt{request} o
\texttt{response}), y el nodo emisor es el par del receptor. En cada par se
comprueba que $L(\text{recepción}) > L(\text{envío})$.

Para analizar la concurrencia, el script reconstruye la relación
\emph{ocurrió antes} ($\rightarrow$) definida por Lamport~\cite{lamport1978}:
dentro de un nodo, cada evento ocurre antes que el siguiente, y el envío de un
mensaje ocurre antes que su recepción. Dos eventos $a$ y $b$ son concurrentes
($a \parallel b$) si no se cumple $a \rightarrow b$ ni $b \rightarrow a$.

\subsection{Resultados}
La Tabla~\ref{tab:lamport} resume las 40 corridas medidas de la sesión
oficial. La columna \emph{Mensajes} cuenta por separado la solicitud y la
respuesta de cada envío, por lo que cada corrida de 100 envíos aporta 200
mensajes. En todas las combinaciones, el reloj de recepción fue mayor que el
de envío en todos los mensajes, y no quedó ningún evento sin pareja. El script
también verifica las corridas de calentamiento y las pruebas previas; sus
resultados quedan en \texttt{analysis/resultados/lamport\_por\_archivo.csv}.

\begin{table}[ht]
\centering
\caption{Verificación de Lamport en la sesión oficial (generada por el script).
\emph{Iguales y conc.}: pares de eventos de nodos distintos con el mismo reloj
que son concurrentes.}
\label{tab:lamport}
\resizebox{\textwidth}{!}{\begin{tabular}{llrrrrr}
\toprule
Conjunto & Mecanismo & Eventos & Mensajes & Cumplen & Fallan & Iguales conc. \\
\midrule
demo\_grpc & gRPC & 60 & 30 & 30 & 0 & 2 \\
tcp\_20\_100b\_eventos & TCP & 80 & 40 & 40 & 0 & 19 \\
tcp\_20\_10kib\_eventos & TCP & 80 & 40 & 40 & 0 & 19 \\
tcp\_concurrente\_10kib\_eventos & TCP & 84 & 42 & 42 & 0 & 68 \\
\bottomrule
\end{tabular}
}
\end{table}

\subsection{Eventos concurrentes con el mismo reloj}
La Tabla~\ref{tab:concurrentes} muestra un ejemplo real por mecanismo y modo.
En el modo concurrente, nodo-A y nodo-B envían su primera solicitud con
$L=1$: ninguno había recibido nada, así que ninguno pudo influir en el otro.
En el modo secuencial, nodo-A recibe la respuesta del mensaje 0001 con $L=4$
al mismo valor con que nodo-B envía la solicitud 0002, y tampoco existe una
cadena de mensajes que conecte esos dos eventos.

\begin{table}[ht]
\centering
\caption{Ejemplos de eventos concurrentes con igual reloj (generada por el script).}
\label{tab:concurrentes}
\resizebox{\textwidth}{!}{\begin{tabular}{lrll}
\toprule
Conjunto & Reloj & Evento 1 & Evento 2 \\
\midrule
demo\_grpc & 20 & nodo-A receive response 0005 & nodo-B send request 0001 \\
tcp\_20\_100b\_eventos & 4 & nodo-A receive response 0001 & nodo-B send request 0002 \\
tcp\_20\_10kib\_eventos & 4 & nodo-A receive response 0001 & nodo-B send request 0002 \\
tcp\_concurrente\_10kib\_eventos & 1 & nodo-A send request 0001 & nodo-B send request 0002 \\
\bottomrule
\end{tabular}
}
\end{table}

\subsection{Por qué Lamport no distingue este caso}
El reloj de Lamport solo garantiza la condición
$a \rightarrow b \Rightarrow L(a) < L(b)$; la implicación inversa no se cumple.
Por eso, dos relojes iguales indican que los eventos no tienen relación causal,
pero un reloj menor no indica que un evento haya causado al otro. Los datos lo
confirman: el script encontró pares de eventos con relojes distintos que también
son concurrentes (columna \texttt{reloj\_distinto\_concurrentes} de
\texttt{lamport\_resumen.csv}). Con un solo número por evento no se puede saber
si dos eventos son concurrentes; para eso se necesitan relojes vectoriales, que
guardan un contador por cada nodo.
